%% notes-tex/genesis/genesis-proposal-redacted/sections/methods.tex
%% [[genesis-proposal-redacted]]
%% https://github.com/Mjvolk3/torchcell/tree/main/notes-tex/genesis/genesis-proposal-redacted/sections/methods.tex
%% SOURCE: hand-transcribed from the proposal docx, "Proposed Research and
%% Methods". Group names replaced by institutional roles; the submission's em-dashes
%% in Aim 2 become commas. Figure 2 is the image embedded in the docx
%% (word/media/image2.png), converted to PDF unchanged.

\section{Proposed research and methods\secstatus{todo}}\label{sec:methods}

\subsection{Aim 1. Extend TorchCell for generalizable foundational models in
metabolic engineering\secstatus{todo}}\label{sec:aim1}

To expand TorchCell to other important industrial organisms, we will start with the
well-studied bacterial strain \org{Escherichia coli} and expand to 10 additional
gram-negative bacteria and 10 additional yeasts. We will then demonstrate
applications on \org{Pseudomonas putida}, which has recently been optimized for
isoprenol production\textsuperscript{11}, an important aviation fuel precursor. In
addition, this case study will help evaluate recommendations provided by TorchCell in
DBTL loops (Fig.~\ref{fig:two}a)\textsuperscript{11}.

Although large foundation models have shown some promising results in mammalian
cells, systematic benchmarking indicates that transformer-based perturbation models
do not consistently outperform simple linear baselines\textsuperscript{13}. This gap
can be filled by better mechanistic representations grounded in physical
constraints, as demonstrated by Merzbacher et al. in their application of deep
learning to flux cone geometry to predict metabolite
production\textsuperscript{14}. Since TorchCell already facilitates the inclusion of
metabolism, we intend to further elaborate its description to ensure deep learning
cell representations are well grounded in physical constraints
(Fig.~\ref{fig:two}c). \emph{Successful applications in organisms with
comparatively less data would demonstrate the successful implementation of MEFMs.}
(Underlined in the submission; italic here because an underline cannot break
across lines.)

\tcfigfit[0.62]{figures/genesis-fig2.pdf}{The figure as embedded in the submission,
which carries no caption of its own. \textbf{a}, organisms entering TorchCell.
\textbf{b}, the perturb-seq style atlas workflow: individual growth of a new yeast
knockout collection, pooling under control and stress conditions, scRNA-seq, and
the genotype-clone-transcriptome readout. \textbf{c}, the metabolic grounding
layers: the metabolic network, the flux cone under mass balance constraints, and
dynamics under kinetic and thermodynamic constraints.}{fig:two}

\subsection{Aim 2. Integrate an improved genome LLM based on GenSLM for metabolic
engineering (HiGenLM) with TorchCell\secstatus{todo}}\label{sec:aim2}

GenSLM was one of the first codon-based LLM models to enable genome-scale predictions
and to predict novel variants in the newly emerging coronavirus genomes
(SARS-CoV-2)\textsuperscript{6}. We are now developing a phenotypically informed DNA
LLM that leverages GenSLM and performs post-training with phenotypic data to make
effective phenotype predictions. Our new version of GenSLM, HiGenLM, focuses on
automatically learning the context of nucleotides, genes, contigs, and operons in
the genome, using a nested, multi-resolution transformer architecture pre-trained on
genomic sequences from thousands of microbial and eukaryotic genomes. We hypothesize
that this hierarchical context learning is critical to add phenotypic awareness. At
its core, HiGenLM decomposes the genome into a hierarchy of overlapping biological
units, codons, genes, operons, regulatory modules, and chromosome segments, and
learns representations at each scale through a combination of masked language
modeling and cross-scale contrastive objectives. The result is a set of scale-aware
genomic embeddings that capture both the fine-grained biochemical specificity of
individual protein-coding sequences and the broader genomic context in which they
operate. Critically, because HiGenLM is pre-trained across phylogenetically diverse
organisms, its representations encode evolutionary conservation signals that
genome-scale models entirely lack: residues under purifying selection, syntenic gene
neighborhoods, and cross-species regulatory grammar are all implicitly encoded in
the embedding geometry. Currently, we are training HiGenLM with phenotypic data by
leveraging Directed Preference Optimization (DPO)\textsuperscript{15}, which will
help us tackle the disparity of phenotypic and genotypic data availability. DPO
works by directly updating the LLM using a loss function that favors preferred
datasets (e.g., higher production phenotypes). An important part of this process
will be augmenting phenotypic data to meet post-training needs, which will be done
by extrapolating from current data and by supplementing phenotypic data from the
TorchCell database.

We will attempt to integrate HiGenLM with TorchCell. HiGenLM gene-level embeddings
can replace or augment the sequence features currently used in TorchCell's
supervised models for fitness, gene interaction, morphology, and expression
prediction. By providing each gene node in TorchCell's graph database with a
HiGenLM-derived embedding giving TorchCell CGT models access to evolutionary and
functional priors that sharply reduce the effective sample complexity for
multi-phenotype prediction. Preliminary evidence from analogous integrations using
GenSLM applied to \org{E. coli} pathway engineering suggests that this embedding
transfer can improve phenotype prediction accuracy by 15--30\% in low-data regimes,
precisely the regime most relevant to metabolic engineering strain development.
Notably, HiGenLM's cross-species pre-training provides a shared embedding space in
which gene function and regulatory logic are transferable across phylogenetic
distance, enabling TorchCell-style phenotype models to be rapidly adapted to new
hosts by grounding their gene representations in HiGenLM's universal genomic
language.

\subsection{Aim 3. Use biofoundries to validate the predictions made by the
proposed AI models and generate large amounts of data in a high-throughput manner
for finetuning of these models\secstatus{todo}}\label{sec:aim3}

Metabolic engineering relies on genetic manipulation to redirect flux within
organisms, necessitating data collection across perturbations of the genome and
environment to construct effective metabolic engineering foundation models. We
propose to collect perturb-seq style single-cell atlas\textsuperscript{12} for
non-model organisms relevant to industrial biotechnology. Data will be used to
evaluate foundation models built on model organisms (e.g., \org{S. cerevisiae} and
\org{E. coli}) and applied to non-model organisms, and to assess model improvements
using high-throughput perturbation data. Selected strains will also be screened with
metabolically relevant metabolites using EchoMS to link high-throughput sequencing
data to more expensive metabolite quantification using previously mentioned
facilities.

To validate the gene target predictions made by the MEFMs, we will focus on model
organisms \org{S. cerevisiae} and \org{E. coli} for which we have established
automated workflows for gene knockouts and overexpression on our iBioFAB and NSF
iBioFoundry (unpublished data). Building upon our prior
studies\textsuperscript{4, 16--17}, we will select $\beta$-carotene as a target
compound for \org{S. cerevisiae}, and lycopene as a target compound for \org{E.
coli}. These organisms and relevant assays have been well established in the lead
group and will help evaluate MEFM predictions for the transfer of pathways to new
hosts.
